% sections/03-methods.tex

\section{Methods: Energy Functionals for the Reconfiguration Graph}
\label{sec:methods}

We define eight energy functionals on $\R(G, 5)$. Each functional
assigns a real number to a proper 5-coloring $c \colon V \to \{1,
\ldots, 5\}$, or to a specific Kempe chain within $c$. Together, they
provide complementary views of the reconfiguration landscape.

\subsection{The Tetrahedral Color Embedding}\label{sec:tet-embed}

We embed the five colors in $\mathbb{R}^3$ as follows:
\begin{equation}\label{eq:tet-embed}
  \ve{1} = (1,1,1), \quad
  \ve{2} = (1,-1,-1), \quad
  \ve{3} = (-1,1,-1), \quad
  \ve{4} = (-1,-1,1), \quad
  \ve{5} = (0,0,0).
\end{equation}

\begin{proposition}[Tetrahedral geometry]\label{prop:tet-geom}
The vectors $\ve{1}, \ldots, \ve{4}$ form the vertices of a regular
tetrahedron with:
\begin{enumerate}[label=(\alph*)]
  \item Pairwise distance: $\norm{\ve{a} - \ve{b}} = 2\sqrt{2}$ for
    all $a \neq b \in \{1,2,3,4\}$.
  \item Centroid at the origin:
    $\frac{1}{4}(\ve{1} + \ve{2} + \ve{3} + \ve{4}) = \vc{0}$.
  \item Inner product: $\ip{\ve{a}}{\ve{b}} = -1$ for $a \neq b$ and
    $\ip{\ve{a}}{\ve{a}} = 3$.
\end{enumerate}
\end{proposition}

\begin{proof}
Direct computation. For distinct $a, b \in \{1,2,3,4\}$, exactly two of
the three coordinates differ in sign, giving
$\norm{\ve{a} - \ve{b}}^2 = 0 + 4 + 4 = 8$ (up to permutation), so
$\norm{\ve{a} - \ve{b}} = 2\sqrt{2}$. The centroid is
$(1+1-1-1, 1-1+1-1, 1-1-1+1)/4 = \vc{0}$. The inner product follows:
$\ip{\ve{a}}{\ve{b}} = (+1) + (-1) + (-1) = -1$, since exactly
one coordinate agrees in sign and two disagree.
\end{proof}

The fifth color $\ve{5} = \vc{0}$ sits at the centroid, representing a
``defect'' or ``unresolved'' color. A 4-coloring uses only
$\ve{1}, \ldots, \ve{4}$; a 5-coloring has at least one vertex at the
origin.

\begin{proposition}[Kempe swaps as rotations]\label{prop:kempe-rotation}
For distinct $a, b \in \{1,2,3,4\}$, let $R_{ab}$ denote the
$180^\circ$ rotation in $\mathrm{SO}(3)$ about the axis
$\vc{n}_{ab} = \frac{1}{2}(\ve{a} + \ve{b})$ (the midpoint of the
tetrahedral edge from $\ve{a}$ to $\ve{b}$). Then $R_{ab}$ exchanges
$\ve{a} \leftrightarrow \ve{b}$ and simultaneously exchanges
$\ve{c} \leftrightarrow \ve{d}$, where
$\{c,d\} = \{1,2,3,4\} \setminus \{a,b\}$. It fixes
$\ve{5} = \vc{0}$. Since the vertices of a Kempe $(a,b)$-chain carry
only colors $a$ and $b$, the exchange of $c$ and $d$ is vacuous on the
chain, and $R_{ab}$ faithfully models the Kempe swap.
\end{proposition}

\begin{proof}
The axis $\vc{n}_{ab}$ passes through both the midpoint of the
edge $\ve{a}\ve{b}$ and the centroid (origin) of the tetrahedron.
A $180^\circ$ rotation about this axis exchanges $\ve{a}$ and
$\ve{b}$. Because the axis also passes through the midpoint of the
opposite edge $\ve{c}\ve{d}$, the rotation simultaneously exchanges
$\ve{c}$ and $\ve{d}$. Since $\ve{5} = \vc{0}$ lies on every axis
through the origin, it is fixed. For an explicit check, take
$a = 1$, $b = 2$: $\vc{n}_{12} = (1, 0, 0)$, and the rotation
$R_{12} = \operatorname{diag}(1, -1, -1)$ sends
$(1,1,1) \mapsto (1,-1,-1) = \ve{2}$ and
$(-1,1,-1) \mapsto (-1,-1,1) = \ve{4}$. The other four cases
follow by the symmetry of the tetrahedron. On a Kempe $(a,b)$-chain,
only colors $a$ and $b$ appear, so the exchange of $c$ and $d$ is
vacuous.
\end{proof}

The six Kempe swap types $\{(1,2), (1,3), (1,4), (2,3), (2,4),
(3,4)\}$ pair into three complementary pairs — $(1,2)$ with $(3,4)$,
$(1,3)$ with $(2,4)$, and $(1,4)$ with $(2,3)$ — sharing the same
rotation axis. There are thus only three distinct $180^\circ$
rotations, and they generate the Klein four-group
$V_4 \cong \mathbb{Z}_2^2$ inside $\mathrm{SO}(3)$. As permutations
of the color set $\{1,2,3,4\}$, however, the six Kempe swap types
are the six transpositions in the symmetric group $S_4$, which
generate $S_4$ in its entirety. We note that $S_4$ is the full
symmetry group of the tetrahedron (including reflections); its
rotation subgroup is $A_4$.

\subsection{Magic Gem Covariance Energy}\label{sec:mg-energy}

Inspired by the Magic Gem framework~\cite{mathewson2025}, we define a
local covariance energy that measures the imbalance of the tetrahedral
color vectors in each vertex neighborhood.

\begin{definition}[Local Magic Gem energy]\label{def:mg-local}
For a vertex $v$ with coloring $c$ and closed neighborhood
$N[v] = \{v\} \cup \{u : u \sim v\}$, define
\begin{equation}\label{eq:mg-local}
  E_{\mathrm{MG}}(v, c) =
  \left\lVert
    \frac{1}{\abs{N[v]}} \sum_{u \in N[v]} \ve{c(u)}
  \right\rVert^2.
\end{equation}
\end{definition}

\begin{definition}[Global Magic Gem energy]\label{def:mg-global}
\begin{equation}\label{eq:mg-global}
  E_{\mathrm{MG}}(G, c) = \sum_{v \in V} E_{\mathrm{MG}}(v, c).
\end{equation}
\end{definition}

\begin{remark}\label{rem:mg-interp}
$E_{\mathrm{MG}}(v, c) = 0$ if and only if the color vectors in the
closed neighborhood of $v$ are ``balanced,'' i.e., their centroid is
the origin. For a vertex of degree 3 colored with all four colors in
its closed neighborhood, this is exactly satisfied. A vertex with
color-5 neighbors contributes zero vectors to the
centroid sum, which may increase or decrease $E_{\mathrm{MG}}$
depending on the existing color distribution. Thus $E_{\mathrm{MG}}$
penalizes both low chromatic diversity
(many neighbors sharing the same color) and the presence of the fifth
color.
\end{remark}

\subsection{Anti-Ferromagnetic Potts Hamiltonian}\label{sec:potts}

\begin{definition}[Extended Potts Hamiltonian]\label{def:potts-ham}
\begin{equation}\label{eq:potts-ham}
  H(c) = J \sum_{(u,v) \in E} \delta_{c(u), c(v)}
         + h \sum_{v \in V} \delta_{c(v), 5},
\end{equation}
where $J > 0$ is the anti-ferromagnetic coupling and $h > 0$ is the
fifth-color penalty field. (We note the sign convention differs from
Section~\ref{sec:phys-models}, which uses $H = -J\sum\delta$ with
$J < 0$; the two conventions agree: both penalize monochromatic edges.)
\end{definition}

For proper colorings, the first term vanishes identically ($J$-term
$= 0$), so the Hamiltonian reduces to
$H(c) = h \cdot \abs{\{v : c(v) = 5\}}$. The energy simply counts
defects (color-5 vertices), weighted by $h$. A 4-coloring has $H = 0$.

\begin{remark}\label{rem:potts-chromatic}
For the standard $q$-state anti-ferromagnetic Potts model (without
the $h$-term), the zero-temperature partition function recovers the
chromatic polynomial:
\begin{equation}
  P(G, q) = \lim_{\beta \to \infty} Z(G, q, \beta),
    \qquad
  Z = \sum_c e^{-\beta J \sum_{E} \delta_{c(u),c(v)}}.
\end{equation}
With the fifth-color penalty ($h > 0$), the zero-temperature limit
of the extended partition function counts proper 4-colorings (only
colors $1, \ldots, 4$), not the full chromatic polynomial in $q$.
At finite temperature, $Z$ weights ``near-colorings'' by their defect
density.
\end{remark}

\subsection{Electrostatic Charge Flow}\label{sec:electrostatic}

We model the fifth color as a charge and study its electrostatic
potential on the graph.

\begin{definition}[Charge density and potential]\label{def:charge}
Define the charge density $\rho \colon V \to \{0, 1\}$ by
$\rho(v) = \delta_{c(v), 5}$. Let $L = D - A$ be the graph Laplacian,
where $D$ is the degree matrix and $A$ is the adjacency matrix. The
electrostatic potential is
\begin{equation}\label{eq:potential}
  \phi = L^+ \rho,
\end{equation}
where $L^+$ denotes the Moore-Penrose pseudoinverse.
\end{definition}

\begin{definition}[Discrete electric field magnitude]\label{def:e-field}
The maximum discrete electric field magnitude at vertex $v$ is
\begin{equation}\label{eq:e-field}
  \mathcal{E}(v) = \max_{u \sim v} \abs{\phi(v) - \phi(u)},
\end{equation}
where $\phi = L^+ \rho$ is the electrostatic potential from
Definition~\ref{def:charge}. (We write $\phi$ for the potential and
$\mathcal{E}$ for the field to avoid notational clash with the vertex
set $V$ and the energy functionals $E$.)
\end{definition}

The potential $\phi$ diffuses the charge density across the graph via
the Laplacian, giving a smooth measure of how ``close'' each vertex is
to color-5 defects. High $\mathcal{E}(v)$ indicates that $v$ sits at
the boundary between a defect-rich and defect-poor region.

\subsection{Surface Tension}\label{sec:surface-tension}

Surface tension quantifies the boundary cost of a Kempe chain.

\begin{definition}[Surface tension of a Kempe chain]\label{def:surface-tension}
For a Kempe $(a,b)$-chain $K$ in coloring $c$, the \emph{surface
tension} is
\begin{equation}\label{eq:surface-tension}
  \sigma(K) = \abs{\{(u, v) \in E : u \in K,\; v \notin K\}},
\end{equation}
the number of boundary edges. The \emph{normalized surface tension} is
\begin{equation}\label{eq:normalized-st}
  \bar{\sigma}(K) = \frac{\sigma(K)}{\abs{K}}.
\end{equation}
\end{definition}

\begin{definition}[Surface tension rigidity]\label{def:st-rigidity}
For a coloring $c$ and color pair $(a,b)$, let
$K_1, \ldots, K_m$ be the Kempe $(a,b)$-chains. The \emph{surface
tension rigidity} is
\begin{equation}\label{eq:st-rigidity}
  \Gamma_{ab}(c) =
  \operatorname{Var}\bigl(\bar{\sigma}(K_1), \ldots,
  \bar{\sigma}(K_m)\bigr),
\end{equation}
where $\operatorname{Var}$ denotes the population variance. (We use
$\Gamma$ rather than $\rho$ to avoid conflict with the charge density
of Definition~\ref{def:charge}.) For $m = 1$ (a single chain),
$\Gamma_{ab}(c) = 0$ by convention.
\end{definition}

\begin{remark}\label{rem:st-interpretation}
$\Gamma_{ab}(c) = 0$ means all chains in the $(a,b)$-subgraph have the
same normalized boundary cost — they are ``equally rigid.'' We will see
(Section~\ref{sec:results}) that this rigidity is a signature of
merge-prone configurations.
\end{remark}

\subsection{Local Entropy}\label{sec:entropy}

\begin{definition}[Local chromatic entropy]\label{def:entropy}
For a vertex $v$ with $\deg(v) = d$, define the local entropy
\begin{equation}\label{eq:entropy}
  S(v, c) = -\sum_{a=1}^{5} p_a \log p_a,
  \qquad
  p_a = \frac{\abs{\{u \sim v : c(u) = a\}}}{\deg(v)},
\end{equation}
with the convention $0 \log 0 = 0$. The global entropy is
$S(G, c) = \sum_v S(v, c)$.
\end{definition}

High local entropy at $v$ means the colors of $v$'s neighbors are
uniformly distributed; low entropy means they are concentrated on one
or two colors. We note that this definition uses the open neighborhood
$N(v)$ (unlike the Magic Gem energy, which uses the closed
neighborhood $N[v]$), and that $p_{c(v)} = 0$ always holds for proper
colorings, so the effective sum has at most $k - 1$ nonzero terms.
We will see that unsafe colorings tend to have higher local entropy
near the critical vertex.

\subsection{Chain Ruggedness (Activation Energy)}
\label{sec:ruggedness}

\begin{definition}[Chain ruggedness]\label{def:ruggedness}
The \emph{ruggedness} (or activation energy) of a Kempe chain $K$ is
\begin{equation}\label{eq:ruggedness}
  \mathcal{A}(K) = \frac{\sigma(K)}{\abs{K}^{2/3}},
\end{equation}
where $\sigma(K)$ is the surface tension
(Definition~\ref{def:surface-tension}). We caution that the exponent
$2/3$ is borrowed from the classical isoperimetric inequality in
$\mathbb{R}^3$ ($S \geq C \cdot V^{2/3}$) and has no rigorous
justification for general graphs: for planar graphs, the discrete
isoperimetric exponent is $1/2$, and for arbitrary graphs it depends
on the expansion properties. The choice of $2/3$ is a modeling
decision motivated by the three-dimensional tetrahedral embedding
of the color space, not by the graph's intrinsic geometry.
\end{definition}

By analogy with protein folding~\cite{levinthal1969}, a high activation
energy corresponds to a large barrier that must be overcome to
reconfigure the chain. Low ruggedness chains are ``floppy'' and easily
swapped; high ruggedness chains resist reconfiguration.

\subsection{Defect Interaction Energy}\label{sec:defect-energy}

\begin{definition}[Defect interaction]\label{def:defect-interaction}
For a coloring $c$ and decay parameter $\lambda > 0$, the defect
interaction energy is
\begin{equation}\label{eq:defect-interaction}
  E_{\text{defect}}(c) =
  \sum_{\substack{\{v, w\} \subseteq V,\; v \neq w \\ c(v) = c(w) = 5}}
  e^{-\lambda \cdot d_G(v, w)},
\end{equation}
where the sum is over unordered pairs of distinct vertices and
$d_G(v, w)$ is the graph distance.
\end{definition}

This Yukawa-type interaction penalizes clustered defects.
For proper colorings, no two color-5 vertices are adjacent (that would
violate properness), so the minimum pairwise distance is $2$ and the
strongest interaction is $e^{-2\lambda}$.
Concentrated defects create localized stress that may force unsafe
swaps during reconfiguration.

\subsection{Computational Pipeline}\label{sec:pipeline}

All energy functionals were computed on the complete set of proper
5-colorings of every planar triangulation on $n \leq 9$ vertices.
Triangulations were generated using the plantri canonical construction
path method (Whitney's theorem guarantees a unique planar embedding for
3-connected graphs), stored as adjacency lists, and verified against
known counts (OEIS A000109). Proper colorings were enumerated
via exhaustive backtracking with symmetry reduction. Kempe chains,
reconfiguration graphs, and BFS distances were computed using standard
graph algorithms implemented in Python with NetworkX. The full
enumeration at $n = 9$ (50 triangulations, ${\sim}282{,}300$
colorings) completes in approximately four minutes on a single core.
Extension to $n = 10$ (233 triangulations, ${\sim}2{,}000{,}000$
colorings) requires approximately two hours.
